\documentclass[11pt,letterpaper]{article}
\usepackage[technical-reference]{cornell-notes}
\usepackage{needspace}

\title{Clause 2: Normative References - Cornell Notes}
\author{}
\date{}
\setDocTitle{Clause 2: Normative References}
\setDocSubtitle{C++ 2024 Cornell Notes}
\setDocOwner{C++ 2024 Cornell Notes}
\setDocAuthor{}
\setDocDate{}
\setCornellCollection{C++ 2024 Cornell Notes}
\setCornellUnitType{Clause}
\setCornellUnitNumber{2}
\setCornellUnitTitle{Normative References}
\setCornellCompanionLabel{C++ Technical Reference}
\hypersetup{pdftitle={Clause 2: Normative References - Cornell Notes},pdfsubject={C++ technical study notes},pdfauthor={}}

\begin{document}
\maketitle
\begin{CornellReferenceOrientationBox}
\textbf{Purpose.} Identifies external technical foundations whose cited content participates in the requirements. The set covers shared vocabulary, mathematical notation, the C language foundation, operating-system interfaces, date and time representation, and a regular-expression grammar.
\par\textbf{Role.} Connects the language and library rules to external definitions without restating them.
\par\textbf{Principal dependencies.} The scope establishes why referenced material matters; later clauses invoke particular external definitions.
\par\textbf{Primary audience.} implementers and readers resolving imported terminology or behavior.
\par\textbf{Major subject groups.} Reference model.
\end{CornellReferenceOrientationBox}
\section{Learning Objectives}
\begin{itemize}[label=\textcolor{CornellReferenceTeal}{$\blacktriangleright$}]
\item Define the governing ideas of Reference model and state the consequences for a program or implementation.
\item Identify the governing ideas of normative reference and state the consequences for a program or implementation.
\item Distinguish the governing ideas of dated reference and state the consequences for a program or implementation.
\item Explain the governing ideas of undated reference and state the consequences for a program or implementation.
\item Trace the governing ideas of imported requirement and state the consequences for a program or implementation.
\end{itemize}
\section{Normative Structure Map}
\small\begin{longtable}{@{}>{\RaggedRight\arraybackslash}p{0.13\textwidth}>{\RaggedRight\arraybackslash}p{0.27\textwidth}>{\RaggedRight\arraybackslash}p{0.19\textwidth}>{\RaggedRight\arraybackslash}p{0.31\textwidth}@{}}
\toprule\textbf{Subclause} & \textbf{Subject} & \textbf{Rule category} & \textbf{Key dependencies / entities}\\\midrule\endfirsthead
\toprule\textbf{Subclause} & \textbf{Subject} & \textbf{Rule category} & \textbf{Key dependencies / entities}\\\midrule\endhead
\bottomrule\endfoot
2.1 & Reference model & core-language semantics & The scope establishes why referenced material matters; later clauses invoke particular external definitions.\\
\end{longtable}\normalsize
\begin{CornellReferenceNormativeBox}A heading maps the clause; it does not replace the detailed conditions. For each operation, read requirements, effects, result, exceptions, complexity, remarks, and notes with their stated normative force.\end{CornellReferenceNormativeBox}
\subsection*{Rule Classification Guide}
\small\begin{longtable}{@{}>{\RaggedRight\arraybackslash}p{0.25\textwidth}>{\RaggedRight\arraybackslash}p{0.67\textwidth}@{}}\toprule\textbf{Label or category} & \textbf{How to read it}\\\midrule\endfirsthead\toprule\textbf{Label or category} & \textbf{How to read it (continued)}\\\midrule\endhead\bottomrule\endfoot
Well-formed / ill-formed & Formation is governed by syntax and semantic rules; an ill-formed program does not become well-formed because one implementation accepts it.\\
Diagnosable rule & A violation requires at least one diagnostic unless the wording places the case outside the diagnosable set.\\
No diagnostic required & The program can be ill-formed while the implementation has no diagnostic obligation.\\
Undefined behavior & No execution requirements are imposed; translation success does not establish safety or portability.\\
Unspecified behavior & One of the permitted outcomes may be chosen without documentation.\\
Implementation-defined behavior & The implementation chooses and documents the behavior.\\
Conditionally supported & The implementation may omit support but documents unsupported constructs.\\
\end{longtable}\normalsize
\section{Cornell Cue-and-Notes}
\Needspace{8\baselineskip}
\subsection*{2.1 Reference model}
\begin{longtable}{@{}>{\RaggedRight\arraybackslash\columncolor{CornellReferenceCueGray}}p{0.28\textwidth}>{\RaggedRight\arraybackslash}p{0.64\textwidth}@{}}
\rowcolor{CornellReferenceNavy}\color{white}\textbf{Cue / recall prompt} & \color{white}\textbf{Notes}\\\endfirsthead
\rowcolor{CornellReferenceNavy}\color{white}\textbf{Cue / recall prompt} & \color{white}\textbf{Notes (continued)}\\\endhead
\arrayrulecolor{CornellReferenceLineGray}
\textbf{What rule applies to Reference model?}\\[-1mm]{\scriptsize\CornellClauseRef{2.1} [intro.\allowbreak{}refs]} & Defines the core-language rules for Reference model, including formation, semantic effect, interactions with type and lookup, and any ill-formed or implementation-dependent cases. Identifies external technical foundations whose cited content participates in the requirements. The set covers shared vocabulary, mathematical notation, the C language foundation, operating-system interfaces, date and time representation, and a regular-expression grammar.\\\hline
\end{longtable}
\begin{CornellReferenceSummaryBox}Reference model supplies the core-language semantics portion of this clause. The key review move is to connect its local wording to the clause-wide model, then classify each condition by normative force before relying on it.\end{CornellReferenceSummaryBox}
\section{Language-Lawyer and Library-Usage Pitfalls}
\begin{CornellReferencePitfallBox}\begin{itemize}[label=\textcolor{CornellReferenceAmber}{$\blacktriangleright$}]
\item Treating every informative mention as a normative dependency.
\item Applying a later edition to a dated reference.
\item Reading an external document without the qualifications imposed by the referring clause.
\item Treating a note, example, or recommended practice as though it imposed a mandatory program rule.
\end{itemize}\end{CornellReferencePitfallBox}
\section{Programmer and Implementation Checklist}
\begin{CornellReferenceChecklistBox}\begin{itemize}[label=$\square$]
\item Determine whether the reference is dated or undated.
\item Read the exact referring context before importing an external rule.
\item Apply any language-specific qualification or adaptation stated by the referring clause.
\item For \CornellClauseRef{2.1}, verify reference model against the precise rule category and all cited dependencies.
\end{itemize}\end{CornellReferenceChecklistBox}
\section{Clause Synthesis}
\begin{CornellReferenceOrientationBox}Identifies external technical foundations whose cited content participates in the requirements. The set covers shared vocabulary, mathematical notation, the C language foundation, operating-system interfaces, date and time representation, and a regular-expression grammar. Connects the language and library rules to external definitions without restating them. A sound review distinguishes syntax and participation from runtime preconditions, keeps notes and examples non-mandatory, and follows cross-clause dependencies before making a portability claim.\end{CornellReferenceOrientationBox}
\section{Self-Test Questions}
\begin{enumerate}
\item What role does Reference model play in the clause's overall model?
\item Which normative distinctions must be kept separate when applying normative reference?
\item How would you audit a program or implementation that depends on dated reference?
\item Scenario: a program relies on an unstated implementation behavior while using undated reference. What must be checked before calling the program portable?
\item Scenario: a program relies on an unstated implementation behavior while using imported requirement. What must be checked before calling the program portable?
\end{enumerate}
\clearpage\section{Answer Key}
\begin{enumerate}
\item Reference model is one part of the clause's core-language semantics model. Its role is understood by connecting the local rule in 2.1 to the clause orientation and its dependent concepts.
\item Keep formation and well-formedness separate from runtime preconditions and effects. Also distinguish mandatory wording from notes, implementation choice from undefined behavior, and interface declarations from detailed contracts; see 2.
\item Audit dated reference by locating the controlling wording in 2, listing inputs and type constraints, proving any preconditions, recording effects and failure behavior, and checking lifetime, invalidation, complexity, and synchronization where applicable.
\item First identify the exact requirement in 2, then classify the relied-on behavior as required, unspecified, implementation-defined, conditionally supported, or undefined. Portability requires satisfying the program obligations and avoiding dependence on a choice not guaranteed in the target environments.
\item First identify the exact requirement in 2, then classify the relied-on behavior as required, unspecified, implementation-defined, conditionally supported, or undefined. Portability requires satisfying the program obligations and avoiding dependence on a choice not guaranteed in the target environments.
\end{enumerate}
\section{Key-Term Recap}
\begin{longtable}{@{}>{\RaggedRight\arraybackslash}p{0.28\textwidth}>{\RaggedRight\arraybackslash}p{0.64\textwidth}@{}}\toprule\textbf{Term} & \textbf{Working meaning}\\\midrule\endfirsthead\toprule\textbf{Term} & \textbf{Working meaning (continued)}\\\midrule\endhead\bottomrule\endfoot
normative reference & An external document whose cited content can form part of a requirement.\\
dated reference & A reference for which the specifically cited edition controls.\\
undated reference & A reference that follows the latest applicable edition, including amendments.\\
imported requirement & A rule made applicable through an explicit normative reference.\\
\end{longtable}
\CornellTakeaway{External material becomes controlling only through the reference model and the precise context that invokes it.}
\end{document}
